\section{Discussion and Conclusion}
\label{sec:discussion}
\subsection{Use in a Multipath Scheduler}
The calibrated upper bound $\hat{q}_{0.975}+c_{h,k}$ is a per-path RTT bound that held about 95\% of the time on unseen sessions. A deadline-aware scheduler can send a packet on a path whose bound at the packet's delivery horizon is below the deadline, and duplicate it on two paths only when no single bound qualifies. Because the DS2S bound follows the input, it is tight on a calm path and loose only around bursts (Fig.~\ref{fig:trace}). A constant-width bound around a mean forecast would declare a path unsafe more often for the same coverage. This is the practical meaning of the lower interval score in Table~\ref{tab:main}. The median forecast suits the ordering of paths by expected delay, where a single spike should not dominate, and the mean suits throughput and capacity estimates, where the squared error is the right loss. Both come from one forward pass.

\subsection{Lessons for Link-State Forecasting}
Three lessons apply beyond our testbed. First, measurement pipelines need a causality check. A time interpolation that is harmless for plots places future samples in the inputs, and its effect is concentrated at the shortest horizons, where schedulers operate. Second, at sub-second resolution the squared error of a point forecast is hard to improve. None of our changes reduced the RTT RMSE by more than 7\% on any link, and a tuned linear model came within one point of the neural baseline's skill on four links. Claims of large point-accuracy gains at this scale should be checked against such references~\cite{basit2024predictability,zeng2023transformers}. Third, the gain of learning is in the shape of the predictive distribution, which only shows when intervals are calibrated on held-out sessions. Calibrating on the test residuals, as is common, gives the nominal coverage by construction and hides the difference.

\subsection{Limitations}
The sessions were recorded at one site, with one ping server and constant-rate iperf traffic, and the link conditions of a moving or airborne terminal may differ. The calibration sessions also serve for model selection and sessions are not exchangeable. The coverage is therefore measured rather than guaranteed. Adaptive conformal updates~\cite{gibbs2021adaptive} can track a drifting offset online. The iperf reports are stamped at the start of their one-second interval, which may let jitter and loss inputs overlap the first second of the horizon for all methods alike. The scheduling benefit is argued from forecast quality, and an end-to-end evaluation inside the MPQUIC client is the next step.

\subsection{Conclusion}
We studied short-horizon forecasting of RTT, jitter, and loss for Starlink, OneWeb, and 5G links at 5\,Hz. We showed that interpolation-based alignment leaks future samples into the inputs and quantified the effect. We proposed an anchored GRU encoder-decoder with non-crossing distributional heads, trained with a squared error plus pinball loss and calibrated across sessions with CQR. On six links and three seeds it improved the 95\% interval score of the RTT forecasts by \ChgMISWorst\% to \ChgMISBest\% and the MAE by \ChgMAEWorst\% to \ChgMAEBest\% against the same network with a scalar output, at a measured coverage of \PICPMin{} to \PICPMax{} and with an unchanged RMSE. Future work will close the loop with the multipath scheduler, add the state of the other paths as inputs, and adapt the calibration online.
